\section{Premis dan Hipotesis}

Hipotesis penelitian ini disusun secara deduktif dari sejumlah premis yang berakar pada Latar Belakang Masalah dan Rumusan Masalah.
Premis-premis tersebut adalah sebagai berikut:

\begin{enumerate}
    \item \textbf{Premis 1:} Perangkat \textit{edge-IoT} memiliki keterbatasan memori, daya komputasi, dan energi, sehingga mekanisme keamanan perlu dirancang ringan agar dapat dijalankan langsung pada perangkat tersebut \parencite{canavese_security_2024}.
    \item \textbf{Premis 2:} Algoritma \textit{compact} seperti \textit{compact-GA} dan \textit{compact-PSO}\label{first:compact-pso} menghemat memori dengan mengasumsikan independensi antar-atribut, sedangkan BOA dan iBOA\label{first:iboa} dapat memodelkan dependensi tetapi tidak dirancang dengan batasan memori yang ketat. Dalam EDA\label{first:eda}, pemodelan dependensi antar-variabel diketahui dapat meningkatkan kualitas solusi pada masalah dengan korelasi antar-variabel yang tinggi dibandingkan model dengan asumsi independensi penuh.
    \item \textbf{Premis 3:} Lingkungan IoT bersifat dinamis, berskala besar, dan memiliki konteks yang terus berubah \parencite{ragothaman_access_2023}, sehingga kebijakan ABAC perlu diperbarui secara berkala. Namun, pembelajaran ulang struktur dependensi pada setiap perubahan data seperti pada pendekatan BOA/iBOA memerlukan biaya komputasi yang tinggi sehingga kurang efisien untuk perangkat berdaya rendah.
    \item \textbf{Premis 4:} Kualitas hasil \textit{policy mining} dipengaruhi oleh kemampuan model probabilistik merepresentasikan dependensi dalam data akses. Oleh karena itu, algoritma \textit{compact} yang tetap memodelkan dependensi antar-atribut diperkirakan dapat menghasilkan kualitas kebijakan yang lebih mendekati algoritma \textit{population-based} seperti ACO dan UBK dibandingkan algoritma \textit{compact} yang mengasumsikan independensi penuh.
\end{enumerate}

Berdasarkan keempat premis tersebut, hipotesis penelitian yang diajukan adalah sebagai berikut:

\begin{enumerate}
    \item \textbf{H1}\label{first:h1} -- CoDA-EDA diharapkan dapat menghasilkan kualitas kebijakan (\textit{F-score}) yang kompetitif dibandingkan dengan keenam algoritma pembanding lainnya (ACO, UBK, BOA, iBOA, \textit{compact-GA, compact-PSO}).
    \item \textbf{H2}\label{first:h2} -- Jejak memori CoDA-EDA tetap datar terhadap ukuran populasi virtual (NP), sehingga diperkirakan secara signifikan lebih rendah dibanding algoritma berpopulasi pada NP besar.
    \item \textbf{H3}\label{first:h3} -- Pemisahan pembaruan struktur dari pembaruan parameter diharapkan dapat mengurangi total waktu komputasi dengan kualitas kebijakan yang non-inferior.
    \item \textbf{H4}\label{first:h4} -- Pada skenario \textit{drift} rendah hingga sedang, CoDA-EDA dengan \textit{warm-start} diharapkan dapat mencapai kualitas kebijakan non-inferior terhadap \textit{re-mining} penuh dalam waktu re-optimasi yang jauh lebih singkat.
\end{enumerate}